\chapter{Linear Programming}
\label{ch28:lp}

In Chapter~27, we studied matrix operations and solved systems of linear equations ($Ax = b$), finding exact balances when constraints are rigid equalities. But real-world resources are rarely exact equalities: you have at most twelve hours of labor, at least five tons of raw material, and a goal to maximize profit. Replacing equalities with linear inequalities and seeking the best possible solution turns linear algebra into \emph{linear programming}.

Linear programming is best understood as a \emph{language}: a single, expressive way to write down an enormous family of optimization problems, paired with general solvers that can answer any problem posed in that language. Once you know the language, shortest paths, maximum flow, and many real scheduling, blending, and resource-allocation problems can all be written as instances of the same geometric problem.

The phrase ``linear programming'' is, like ``dynamic programming,'' a
historical accident. ``Programming'' here means \emph{planning} -- arranging a
schedule or an allocation -- not writing code. And ``linear'' means exactly
what it says: every quantity we care about is built only out of additions and
constant multiples. No squares, no products of two unknowns, no \texttt{if}
branches. Restricting every constraint and objective to straight lines and flat
hyperplanes is what makes general, efficient solvers possible.

\begin{intuitionbox}
Picture a factory that makes two products. Each product earns some profit per
unit, and each consumes some of your limited labor hours and raw material. You
want to choose how much of each to make to maximize profit, without using more
labor or material than you have. Every constraint (``labor used $\le$ labor
available'') is a straight-line boundary cutting through the plane of possible
production plans. Together those boundaries fence off a region of \emph{legal}
plans, and ``maximize profit'' means sliding a ruler across that region until
it just touches the farthest corner. The rest of the chapter makes this picture
precise and shows how to find the best corner quickly.
\end{intuitionbox}

We begin with the algebra and geometry of two-variable programs, trace how the simplex algorithm walks along the edges of multi-dimensional polytopes, uncover the duality theorem that certifies optimality, and conclude by using linear programming to approximate hard discrete problems.

\section{Modeling: From Words to Algebra}
\label{ch28:modeling}

Let us make the factory concrete. You run a small workshop that builds two
kinds of furniture, chairs and tables. The numbers:

\begin{center}
\begin{tabular}{l rr r}
\toprule
 & chair & table & available \\
\midrule
profit per unit (\$) & 3 & 5 & --- \\
wood (board-feet)    & 1 & 2 & 8 \\
labor (hours)        & 2 & 2 & 12 \\
\bottomrule
\end{tabular}
\end{center}

Each chair earns \$3 and each table earns \$5. A chair needs 1 board-foot of
wood and 2 hours of labor; a table needs 2 board-feet and 2 hours. You have 8
board-feet of wood and 12 hours of labor this week. How many chairs and tables
should you build to maximize profit?

The first move in every modeling problem is to name the things you get to
\emph{decide}. Call them \emph{decision variables}. Let $x_1$ be the number of
chairs and $x_2$ the number of tables. Now everything else is forced.

\begin{itemize}
\item \textbf{Objective.} Total profit is $3x_1 + 5x_2$, and you want it as
  large as possible. A linear function of the variables that you push up or
  down is called the \emph{objective function}.
\item \textbf{Constraints.} You cannot use more wood than you have:
  $x_1 + 2x_2 \le 8$. You cannot use more labor than you have:
  $2x_1 + 2x_2 \le 12$. And you cannot build a negative number of anything:
  $x_1 \ge 0$ and $x_2 \ge 0$. Those last two are the \emph{nonnegativity
  constraints}, and they appear so often that we usually group them on a single
  line.
\end{itemize}

Put it all together and you have a \emph{linear program}:
\[
\begin{aligned}
\text{maximize}\quad   & 3x_1 + 5x_2 \\
\text{subject to}\quad & x_1 + 2x_2 \le 8 \\
                       & 2x_1 + 2x_2 \le 12 \\
                       & x_1, x_2 \ge 0.
\end{aligned}
\]

\begin{focusbox}
An LP has exactly three ingredients, and the word ``linear'' applies to all of
them:
\begin{enumerate}
\item \textbf{Decision variables} $x_1,\dots,x_n$ -- the quantities you choose.
\item A \textbf{linear objective} $c_1 x_1 + \cdots + c_n x_n$ to maximize (or
  minimize).
\item \textbf{Linear constraints}, each an inequality ($\le$ or $\ge$) or
  equality ($=$) between a linear combination of the variables and a constant.
\end{enumerate}
``Linear'' means every appearance of a variable is just that variable times a
constant, added to others. No $x_1^2$, no $x_1 x_2$, no $\sqrt{x_1}$.
\end{focusbox}

\begin{trapbox}
The single most common modeling mistake is sneaking in a \emph{nonlinearity}
without noticing. ``The profit per chair drops once I sell more than 10'' is a
nonlinear objective. ``I can build chairs \emph{or} tables but not both'' is a
logical \texttt{or}, which is also not linear (it needs integer 0/1 variables --
see Section~\ref{ch28:relaxation}). Watch for products of variables (``revenue
$=$ price $\times$ quantity'' when \emph{both} are decisions) -- that is the
classic trap. If your model has any of these, it is not a plain LP, and the
gentle theory of this chapter does not directly apply.
\end{trapbox}

\begin{aestheticbox}
This template fits a very wide range of real decisions. Oil
refineries blend crude streams into gasoline grades to hit octane targets at
least cost -- an LP. Airlines assign crews and aircraft to routes -- LPs (with
integers) at industrial scale. Hospitals schedule nurses, farms choose crop
mixes, data centers route traffic, portfolios balance risk. A
seventy-year-old method is still in constant use because ``maximize a linear
thing subject to linear limits'' describes the basic structure of a great many
planning problems.
\end{aestheticbox}

\subsection{A second model: shortest paths as an LP}

To convince you LP is a \emph{language}, let us encode something you already
know -- single-source shortest paths -- as a linear program. Suppose you want
the shortest-path distances $d_v$ from a source $s$ in a graph with edge
weights $w(u,v)$. The trick is to ask for the distances that are as
\emph{large} as possible while still respecting the triangle inequality on
every edge.
\[
\begin{aligned}
\text{maximize}\quad   & \sum_{v} d_v \\
\text{subject to}\quad & d_v - d_u \le w(u,v) \quad \text{for every edge }(u,v),\\
                       & d_s = 0.
\end{aligned}
\]
The constraint $d_v \le d_u + w(u,v)$ is exactly the relaxation condition from
the shortest-paths chapter. Maximizing the total pushes each $d_v$ up until
some constraint stops it -- and the value where it gets stuck is precisely the
true shortest-path distance. We will not prove that here, but file it away: a
problem you solved with Dijkstra is \emph{also} an LP. This generality is what
makes LP useful.

\section{Standard and Slack Forms}
\label{ch28:forms}

If we want general algorithms, ``any LP'' has to mean something precise.
Linear programs can be written in many surface forms -- maximizing or
minimizing, mixing $\le$, $\ge$, and $=$ constraints, or allowing unconstrained
variables. Two canonical forms standardize all of them.

\subsection{Standard form}

\begin{formalbox}
A linear program is in \textbf{standard form} if it is a \emph{maximization}
over nonnegative variables subject to $\le$ constraints:
\[
\begin{aligned}
\text{maximize}\quad   & \sum_{j=1}^{n} c_j x_j \\
\text{subject to}\quad & \sum_{j=1}^{n} a_{ij} x_j \le b_i \quad (i = 1,\dots,m),\\
                       & x_j \ge 0 \quad (j = 1,\dots,n).
\end{aligned}
\]
In matrix shorthand, with $c,x \in \R^n$, $b \in \R^m$, and $A$ an $m\times n$
matrix: \emph{maximize $c^\top x$ subject to $Ax \le b$ and $x \ge 0$.}
\end{formalbox}

Our furniture problem is already in standard form, with
\[
c = \begin{pmatrix}3\\5\end{pmatrix},\quad
A = \begin{pmatrix}1 & 2\\ 2 & 2\end{pmatrix},\quad
b = \begin{pmatrix}8\\12\end{pmatrix}.
\]

Every LP can be massaged into standard form with a few mechanical moves. Keep
this little phrasebook:

\begin{itemize}
\item \textbf{Minimize $\to$ maximize.} To minimize $c^\top x$, instead
  maximize $-c^\top x$ and negate the answer at the end.
\item \textbf{$\ge$ constraint $\to$ $\le$.} Multiply both sides by $-1$:
  the constraint $\sum_j a_{ij} x_j \ge b_i$ becomes
  $\sum_j (-a_{ij}) x_j \le -b_i$.
\item \textbf{Equality $\to$ two inequalities.} Replace
  $\sum_j a_{ij}x_j = b_i$ by the pair $\le b_i$ and $\ge b_i$.
\item \textbf{Free (unbounded) variable.} If some $x_j$ is allowed to be
  negative, split it: write $x_j = x_j' - x_j''$ with $x_j', x_j'' \ge 0$. Any
  real value is some nonnegative minus some nonnegative.
\end{itemize}

\begin{example}
Convert ``minimize $2x_1 - 3x_2$ subject to $x_1 + x_2 = 4$, $x_1 \ge 1$, with
$x_2$ free'' to standard form. Maximize $-2x_1 + 3x_2$. Split $x_2 = x_2' -
x_2''$. The equality becomes $x_1 + x_2' - x_2'' \le 4$ and
$-x_1 - x_2' + x_2'' \le -4$. The bound $x_1 \ge 1$ becomes $-x_1 \le -1$. Now
every variable ($x_1, x_2', x_2''$) is $\ge 0$ and every constraint is $\le$.
Done.
\end{example}

\subsection{Slack form}

Standard form is the right shape for stating theory and duality. But the
simplex algorithm prefers a cousin in which every inequality has been turned
into an \emph{equality} by introducing a brand-new nonnegative variable that
soaks up the difference. That extra variable is called a \emph{slack
variable}, because it measures how much slack -- unused resource -- remains in
the constraint.

\begin{intuitionbox}
Geometrically, a constraint like $\sum_j a_{ij} x_j \le b_i$ divides space into
two halves separated by the boundary plane $\sum_j a_{ij} x_j = b_i$. If your
current production plan $x$ uses strictly less than $b_i$, you sit strictly
inside the legal half-space. The slack variable
\[
x_{n+i} = b_i - \sum_{j=1}^n a_{ij} x_j
\]
measures your remaining buffer to that boundary wall.
Requiring $x_{n+i} \ge 0$ is mathematically identical to the original
inequality, but it turns the
inequality into a simple linear equation. When $x_{n+i} = 0$, you are standing
directly on the boundary hyperplane -- the constraint is \emph{tight} (or
\emph{active}). When $x_{n+i} > 0$, you have breathing room.
\end{intuitionbox}

To convert a standard-form LP into slack form systematically:
\begin{enumerate}
\item For each constraint $i \in \{1,\dots,m\}$, create a fresh variable
  $x_{n+i}$ and define it as the difference between the bound $b_i$ and the
  left-hand side $\sum_{j=1}^n a_{ij} x_j$.
\item Express each basic variable $x_{n+1},\dots,x_{n+m}$ solely in terms
  of the original (\emph{nonbasic}) variables $x_1,\dots,x_n$.
\item Write the objective $z$ explicitly as an offset $v$ (initially $0$) plus a
  linear combination of the nonbasic variables.
\item Require all variables -- both original and slack -- to be nonnegative:
  $x_k \ge 0$ for all $k \in \{1,\dots,n+m\}$.
\end{enumerate}

Doing this for both constraints of our furniture LP:
\[
\begin{aligned}
z   &= 0 + 3x_1 + 5x_2 \\
x_3 &= 8 - x_1 - 2x_2 \\
x_4 &= 12 - 2x_1 - 2x_2,
\end{aligned}
\qquad x_1,x_2,x_3,x_4 \ge 0 .
\]

\begin{focusbox}
\textbf{Slack form} writes the objective $z$ and each slack variable on the
left, as a function of the others on the right, plus a constant. The variables
that sit alone on the left are called \emph{basic}; the ones on the right are
\emph{nonbasic}. A slack variable being zero means its constraint is
\emph{tight} (the resource is fully used); positive slack means room to spare.
The simplex method keeps track of its progress through this split into basic
and nonbasic variables.
\end{focusbox}

A key geometric consequence: if you set all $n$ nonbasic variables to $0$, you
can read off the $m$ basic variables instantly. Here, setting $x_1 = x_2 = 0$
gives $x_3 = 8$, $x_4 = 12$, and $z = 0$. In $n$-dimensional space, setting $n$
variables to zero pins down the intersection of $n$ bounding hyperplanes -- a
single sharp corner (vertex) of the polytope! Here that vertex is the origin
$(0,0)$, representing the ``build nothing'' plan -- a valid starting point that
we can now systematically climb away from.

\section{The Geometry of an LP}
\label{ch28:geometry}

Before solving the system algebraically, let us plot the furniture constraints
in the plane so we can see the feasible region directly.

Each variable is an axis. With two variables we live in the plane, with $x_1$
(chairs) horizontal and $x_2$ (tables) vertical. Each constraint is a half of
the plane:
\begin{itemize}
\item $x_1 \ge 0$: everything to the right of the vertical axis.
\item $x_2 \ge 0$: everything above the horizontal axis.
\item $x_1 + 2x_2 \le 8$: everything on or below the line through $(8,0)$ and
  $(0,4)$.
\item $2x_1 + 2x_2 \le 12$, i.e. $x_1 + x_2 \le 6$: on or below the line
  through $(6,0)$ and $(0,6)$.
\end{itemize}

The points satisfying \emph{all} the constraints form the shaded region in
Figure~\ref{ch28:fig-polytope}, called the \emph{feasible region}. Because it
is an intersection of half-planes, it is a convex polygon (in higher dimensions,
a convex \emph{polytope}).

\begin{figure}[H]
\centering
\begin{tikzpicture}[scale=0.95]
  % Grid and subtle axes guide
  \draw[step=1, black!8, very thin] (-0.5,-0.5) grid (8.5,6.5);

  % Axes
  \draw[->, thick, black!70] (-0.3,0) -- (8.8,0) node[right, font=\small] {$x_1$ (chairs)};
  \draw[->, thick, black!70] (0,-0.3) -- (0,6.8) node[above, font=\small] {$x_2$ (tables)};

  % Axis ticks (selective to avoid colliding with vertex labels)
  \foreach \x in {2,8} {
    \draw[black!60] (\x,0.08) -- (\x,-0.08) node[below, font=\footnotesize] {$\x$};
  }
  \foreach \y in {2,6} {
    \draw[black!60] (0.08,\y) -- (-0.08,\y) node[left, font=\footnotesize] {$\y$};
  }

  % Constraint half-plane boundaries
  % Wood: x1 + 2x2 = 8 -> (8,0) to (0,4)
  \draw[sourceblue, very thick] (-0.2, 4.1) -- (8.4, -0.2);
  \node[sourceblue!85!black, font=\footnotesize, rotate=-26.6, above, fill=white, inner sep=1.5pt] at (6.4, 1.0) {Wood: $x_1 + 2x_2 \le 8$};

  % Labor: x1 + x2 = 6 -> (6,0) to (0,6)
  \draw[targetgreen, very thick] (-0.3, 6.3) -- (6.3, -0.3);
  \node[targetgreen!85!black, font=\footnotesize, rotate=-45, above right, fill=white, inner sep=1.5pt] at (0.8, 5.4) {Labor: $x_1 + x_2 \le 6$};

  % Feasible region shading
  \fill[sourceblue!15, opacity=0.85] (0,0) -- (6,0) -- (4,2) -- (0,4) -- cycle;
  \draw[black!70, thick] (0,0) -- (6,0) -- (4,2) -- (0,4) -- cycle;
  \node[sourceblue!80!black, font=\small\bfseries] at (2.0, 0.45) {Feasible Region};

  % Iso-profit contour lines: 3x1 + 5x2 = const => x2 = -0.6 x1 + const/5
  % z = 14: sub-optimal interior contour
  \draw[dashed, black!40, thick] (0, 2.8) -- (4.67, 0);
  \node[black!55, font=\tiny, rotate=-31, above, fill=white, inner sep=1pt] at (3.8, 0.52) {$z = 14$};

  % z = 22 (Optimal tangent line): (0, 4.4) to (7.33, 0)
  \draw[dashed, trapred!80!black, very thick] (0, 4.4) -- (7.33, 0);
  \node[trapred!90!black, font=\scriptsize\bfseries, rotate=-31, above, fill=white, inner sep=1.5pt] at (1.1, 3.74) {$z = 22$};

  % Objective gradient vector c = (3, 5) perpendicular to contours
  \draw[->, trapred, line width=1.5pt] (1.0, 0.8) -- ++(0.5, 0.833)
    node[midway, left=2pt, font=\scriptsize\bfseries, trapred, fill=white, inner sep=1pt] {$\nabla z = \begin{psmallmatrix} 3 \\ 5 \end{psmallmatrix}$};

  % Polytope Vertices
  \fill[black!80] (0,0) circle (2.5pt);
  \node[below left, font=\footnotesize, align=right] at (-0.1,-0.1) {$O(0,0)$\\ {\scriptsize\color{black!60}$z=0$}};

  \fill[black!80] (6,0) circle (2.5pt);
  \node[below=4pt, font=\footnotesize, fill=white, inner sep=1pt] at (5.85, -0.04) {$A(6,0)$};

  \fill[black!80] (0,4) circle (2.5pt);
  \node[left, font=\footnotesize, align=right] at (-0.2, 4) {$C(0,4)$\\ {\scriptsize\color{black!60}$z=20$}};

  % Optimum vertex B(4,2) with halo and dedicated upper callout
  \draw[trapred, thick] (4,2) circle (5.5pt);
  \fill[trapred] (4,2) circle (2.8pt);
  \node[above right=6pt, font=\footnotesize\bfseries, trapred!90!black, align=left] at (4, 2.2)
    {$B(4,2)$ \textbf{[Optimal]}\\[1pt]\scriptsize\color{black!70}$z = 3(4) + 5(2) = 22$};
\end{tikzpicture}
\caption{The geometry of the furniture linear program. The feasible region
(shaded blue) is bounded by four half-planes: nonnegativity ($x_1 \ge 0, x_2 \ge 0$),
the wood constraint ($x_1 + 2x_2 \le 8$, blue boundary), and the labor constraint
($x_1 + x_2 \le 6$, green boundary). Dashed lines trace iso-profit contours
$z = 3x_1 + 5x_2 = \text{const}$, advancing perpendicular to the objective gradient
$\vec{c} = (3,5)$. The maximum achievable profit touches the feasible polygon at the
extreme corner $B(4,2)$ with value $z = 22$.}
\label{ch28:fig-polytope}
\end{figure}

The four corners are $O=(0,0)$, $A=(6,0)$, $B=(4,2)$, and $C=(0,4)$. (Corner
$B$ is where the two constraint lines cross; solve $x_1+2x_2=8$ and
$x_1+x_2=6$ simultaneously: subtracting gives $x_2 = 2$, then $x_1 = 4$.)

Now the objective. The function $z = 3x_1 + 5x_2$ has the same value along any
line of slope $-3/5$; these are the dashed \emph{iso-profit} lines. Increasing
$z$ slides that line in the direction of the gradient $(3,5)$ -- up and to the
right. To maximize $z$ you push the iso-profit line as far in the $(3,5)$
direction as you can while still touching the feasible region. The last point
it touches is the optimum.

Let us just evaluate $z$ at all four corners:
\[
z(O)=0,\quad z(A)=18,\quad z(B)=3\cdot4+5\cdot2=22,\quad z(C)=5\cdot4=20.
\]
The winner is $B=(4,2)$ with profit $\boxed{22}$: build 4 chairs and 2 tables.

\begin{intuitionbox}
Notice the optimum landed on a \emph{corner}, not in the middle of an edge or
the interior. This always happens: an iso-profit line, sliding outward, leaves the
region last at a vertex (or, in a tie, along a whole edge -- but then a vertex
of that edge is still optimal). So you never need to search the infinite
interior. \textbf{You only ever need to check the corners.} A continuous
problem with infinitely many feasible points collapses to a finite search over
vertices. The simplex method is built on this observation.
\end{intuitionbox}

\begin{formalbox}
The feasible region of an LP is a convex set: if $x$ and $y$ are both feasible
then so is every point $\lambda x + (1-\lambda)y$ on the segment between them,
for $0 \le \lambda \le 1$. (Each constraint is preserved under averaging
because it is linear.) A \textbf{vertex} (or \emph{basic feasible solution}) is
a feasible point that is \emph{not} the midpoint of any segment lying in the
region -- a true corner. The fundamental theorem of linear programming says: if
an LP has an optimal solution, then some \emph{vertex} is optimal.
\end{formalbox}

\begin{trapbox}
``Just check the corners'' sounds like a finished algorithm, but it hides a
trap: the number of vertices can be \emph{exponential} in the number of
variables and constraints. An LP in $n$ variables with $m$ constraints can have
on the order of $\binom{m+n}{n}$ vertices. Enumerating them all is hopeless for
large problems. What the simplex method adds is a rule for walking from vertex to neighboring
vertex, always uphill, so that in practice you reach the best one after visiting
only a small fraction of them.
\end{trapbox}

A few more geometric vocabulary words you will meet:
\begin{itemize}
\item \textbf{Infeasible.} The constraints contradict each other and the
  feasible region is empty. (E.g. $x \le 1$ and $x \ge 2$.) There is nothing to
  optimize.
\item \textbf{Unbounded.} The region runs off to infinity in a direction where
  the objective keeps improving, so there is no finite maximum. (E.g. maximize
  $x_1$ with only $x_1 \ge 0$.)
\item \textbf{Bounded and feasible.} The good case: a finite optimum exists, on
  a vertex.
\end{itemize}

\section{The Simplex Method, by Example}
\label{ch28:simplex}

Now we make ``walk uphill from corner to corner'' mechanical. The simplex
method, due to George Dantzig in 1947, does exactly this, using slack form as
its ledger. We will run it on the furniture LP and \emph{watch} it travel from
$O$ to $B$.

\subsection{The idea in one paragraph}

Start at a known vertex (the origin, where all original variables are 0 and the
slacks hold all the resource). Look at the objective written in terms of the
current nonbasic variables. If some nonbasic variable has a \emph{positive}
coefficient, increasing it would raise $z$ -- so we want to bring it into the
basis. We increase it as much as we can until some basic variable hits 0 (a
constraint goes tight); that variable leaves the basis. Swap them, rewrite
everything, and repeat. Each swap moves us to an adjacent vertex with no worse
(usually strictly better) objective value. When no nonbasic variable has a
positive coefficient, no local move helps -- and by convexity, we are at the
global optimum.

\begin{focusbox}
Two pieces of jargon for the swap:
\begin{itemize}
\item The variable we increase (positive objective coefficient) is the
  \textbf{entering variable} -- it enters the basis.
\item The basic variable that hits 0 first, limiting how far we can go, is the
  \textbf{leaving variable} -- it leaves the basis.
\end{itemize}
One entering, one leaving, per iteration. The rewrite that re-expresses
everything in terms of the new basis is called a \textbf{pivot}, by analogy
with Gaussian elimination.
\end{focusbox}

\subsection{Iteration 0: the starting dictionary}

We begin in slack form, which when written this way is often called a
\emph{dictionary}:
\[
\begin{aligned}
z   &= \phantom{8-{}}3x_1 + 5x_2 \\
x_3 &= 8 - x_1 - 2x_2 \\
x_4 &= 12 - 2x_1 - 2x_2.
\end{aligned}
\]
Basic variables: $x_3, x_4$. Nonbasic: $x_1, x_2$ (set to 0). Current vertex:
$(x_1,x_2)=(0,0)$, with $z = 0$. This is corner $O$.

\textbf{Choose entering variable.} Both $x_1$ (coefficient $3$) and $x_2$
(coefficient $5$) would increase $z$. A common rule is to pick the largest
coefficient, so let $x_2$ enter.

\textbf{Choose leaving variable (ratio test).} As we raise $x_2$ from 0 (keeping
$x_1=0$), how far can we go before a basic variable goes negative? Look at each
basic equation:
\[
x_3 = 8 - 2x_2 \ge 0 \;\Rightarrow\; x_2 \le \tfrac{8}{2}=4,\qquad
x_4 = 12 - 2x_2 \ge 0 \;\Rightarrow\; x_2 \le \tfrac{12}{2}=6.
\]
The tightest bound wins: $x_2$ can rise to $4$, at which point $x_3$ hits $0$.
So $x_2$ enters and $x_3$ leaves. This is the \emph{ratio test}: divide each
constant by the (positive) coefficient of the entering variable and take the
minimum.

\subsection{Iteration 1: the pivot}

Solve the binding equation $x_3 = 8 - x_1 - 2x_2$ for the entering variable
$x_2$:
\[
x_2 = 4 - \tfrac12 x_1 - \tfrac12 x_3 .
\]
Now substitute this everywhere $x_2$ appeared. In the objective:
\[
z = 3x_1 + 5\bigl(4 - \tfrac12 x_1 - \tfrac12 x_3\bigr)
  = 20 + \tfrac12 x_1 - \tfrac52 x_3 .
\]
In the other basic equation:
\[
x_4 = 12 - 2x_1 - 2\bigl(4 - \tfrac12 x_1 - \tfrac12 x_3\bigr)
    = 4 - x_1 + x_3 .
\]
So the new dictionary is
\[
\begin{aligned}
z   &= 20 + \tfrac12 x_1 - \tfrac52 x_3 \\
x_2 &= 4 - \tfrac12 x_1 - \tfrac12 x_3 \\
x_4 &= 4 - x_1 + x_3 .
\end{aligned}
\]
Basic: $x_2, x_4$. Nonbasic: $x_1, x_3$ (both 0). Vertex:
$(x_1,x_2)=(0,4)$ -- that is corner $C$ -- with $z = 20$. We climbed from 0 to
20 in one step.

\textbf{Are we done?} The objective row is $z = 20 + \tfrac12 x_1 - \tfrac52
x_3$. The coefficient on $x_1$ is $+\tfrac12 > 0$, so raising $x_1$ still helps.
Not done. Let $x_1$ enter.

\textbf{Ratio test for $x_1$:}
\[
x_2 = 4 - \tfrac12 x_1 \ge 0 \Rightarrow x_1 \le \tfrac{4}{1/2}=8,\qquad
x_4 = 4 - x_1 \ge 0 \Rightarrow x_1 \le \tfrac{4}{1}=4 .
\]
Minimum is $4$, achieved by the $x_4$ row, so $x_4$ leaves.

\subsection{Iteration 2: the second pivot}

Solve the binding equation $x_4 = 4 - x_1 + x_3$ for $x_1$:
\[
x_1 = 4 - x_4 + x_3 .
\]
Substitute into the objective:
\[
z = 20 + \tfrac12(4 - x_4 + x_3) - \tfrac52 x_3
  = 22 - 2 x_3 - \tfrac12 x_4 .
\]
And into the $x_2$ row:
\[
x_2 = 4 - \tfrac12(4 - x_4 + x_3) - \tfrac12 x_3
    = 2 - x_3 + \tfrac12 x_4 .
\]
New dictionary:
\[
\begin{aligned}
z   &= 22 - 2x_3 - \tfrac12 x_4 \\
x_1 &= 4 + x_3 - x_4 \\
x_2 &= 2 - x_3 + \tfrac12 x_4 .
\end{aligned}
\]
Basic: $x_1, x_2$. Nonbasic: $x_3, x_4$ (both 0). Vertex:
$(x_1,x_2) = (4,2)$ -- corner $B$ -- with $z = 22$.

\textbf{Stop?} Every coefficient in the objective row is now \emph{negative}
($-2$ and $-\tfrac12$). Increasing either nonbasic variable would only lower
$z$. There is no improving move. \textbf{We are optimal:} build 4 chairs and 2
tables for \$22, exactly matching what the corner-by-corner geometry told us.

Figure~\ref{ch28:fig-walk} draws the sequence of pivots on the feasible polygon:
two pivots, two edges, climbing $0 \to 20 \to 22$, alongside the algebraic ledger
of dictionary states.

\begin{figure}[H]
\centering
\begin{tikzpicture}[scale=0.82, every node/.style={transform shape}]
  % Left panel: Geometry of pivots
  \begin{scope}[shift={(0,0)}]
    \node[anchor=south west, font=\bfseries\small, black!80] at (-0.2, 5.5) {(a) Polytope Vertex Traversal};
    % Axes
    \draw[->, thick, black!60] (-0.2,0) -- (6.8,0) node[right, font=\footnotesize] {$x_1$};
    \draw[->, thick, black!60] (0,-0.2) -- (0,4.9) node[above, font=\footnotesize] {$x_2$};

    % Feasible polygon
    \fill[sourceblue!12] (0,0) -- (6,0) -- (4,2) -- (0,4) -- cycle;
    \draw[black!30, thin] (0,0) -- (6,0) -- (4,2) -- (0,4) -- cycle;

    % Constraint lines
    \draw[sourceblue!60, dashed] (7.2,0.4) -- (0,4);
    \draw[targetgreen!60, dashed] (6,0) -- (1.0, 5.0);

    % Path O -> C (Pivot 1)
    \draw[->, trapred, line width=2.0pt] (0,0) -- (0,3.88);
    \node[trapred!90!black, left, font=\scriptsize\bfseries] at (-0.05, 2.0) {Pivot 1};

    % Path C -> B (Pivot 2)
    \draw[->, trapred, line width=2.0pt] (0,4) -- (3.88,2.06);
    \node[trapred!90!black, above right, font=\scriptsize\bfseries] at (1.8, 3.1) {Pivot 2};

    % Vertices
    % O(0,0)
    \fill[black!80] (0,0) circle (2.5pt);
    \node[below left, font=\footnotesize] at (0,0) {$O(0,0)$};
    \node[below left, font=\tiny, black!70] at (-0.05,-0.3) {$z=0$};

    % C(0,4)
    \fill[black!80] (0,4) circle (2.5pt);
    \node[left, font=\footnotesize] at (-0.12,4.15) {$C(0,4)$};
    \node[left, font=\tiny, black!70] at (-0.12,3.72) {$z=20$};

    % B(4,2)
    \draw[trapred, thick] (4,2) circle (5pt);
    \fill[trapred] (4,2) circle (2.5pt);
    \node[above right=2pt, font=\footnotesize\bfseries, trapred!90!black, fill=white, inner sep=1pt, align=left] at (4.05,2.05) {$B(4,2)$ [Opt]\\[-1pt]{\tiny\normalfont\color{black!70}$z=22$}};

    % Unvisited A(6,0)
    \fill[black!35] (6,0) circle (2pt);
    \node[below, font=\footnotesize, black!50] at (6,-0.06) {$A(6,0)$};
    \node[below, font=\tiny, black!40] at (6,-0.45) {($z=18$, skipped)};
  \end{scope}

  % Right panel: Algebraic Dictionary Ledger
  \begin{scope}[shift={(7.3, 0.15)}]
    \node[anchor=south west, font=\bfseries\small, black!80] at (-0.4, 5.35) {(b) Simplex Dictionary Transitions};

    % Step 0 Card
    \draw[rounded corners=3pt, draw=black!25, fill=black!2, line width=0.6pt] (-0.4, 3.35) rectangle (8.7, 4.85);
    \node[anchor=west, font=\footnotesize\bfseries, trapred!90!black] at (-0.2, 4.55) {Step 0: At $O(0,0)$, $z = 0$};
    \node[anchor=west, font=\scriptsize] at (-0.2, 4.15) {Nonbasic: $\{x_1=0, x_2=0\}$ \quad Basic: $\{x_3=8, x_4=12\}$};
    \node[anchor=west, font=\scriptsize, black!80] at (-0.2, 3.7) {Enter $x_2$ ($+5$) $\xrightarrow{\text{ratio}} \min(8/2, 12/2)=4 \implies$ Leave $x_3$};

    % Step 1 Card
    \draw[rounded corners=3pt, draw=black!25, fill=black!2, line width=0.6pt] (-0.4, 1.65) rectangle (8.7, 3.15);
    \node[anchor=west, font=\footnotesize\bfseries, trapred!90!black] at (-0.2, 2.85) {Step 1: At $C(0,4)$, $z = 20$};
    \node[anchor=west, font=\scriptsize] at (-0.2, 2.45) {Nonbasic: $\{x_1=0, x_3=0\}$ \quad Basic: $\{x_2=4, x_4=4\}$};
    \node[anchor=west, font=\scriptsize, black!80] at (-0.2, 2.0) {Enter $x_1$ ($+\tfrac12$) $\xrightarrow{\text{ratio}} \min(4/\tfrac12, 4/1)=4 \implies$ Leave $x_4$};

    % Step 2 Card (Optimal)
    \draw[rounded corners=3pt, draw=targetgreen!60!black, fill=targetgreen!6, line width=0.8pt] (-0.4, -0.05) rectangle (8.7, 1.45);
    \node[anchor=west, font=\footnotesize\bfseries, targetgreen!80!black] at (-0.2, 1.15) {Step 2: At $B(4,2)$, $z = 22$ [Optimal]};
    \node[anchor=west, font=\scriptsize] at (-0.2, 0.75) {Nonbasic: $\{x_3=0, x_4=0\}$ \quad Basic: $\{x_1=4, x_2=2\}$};
    \node[anchor=west, font=\scriptsize\bfseries, targetgreen!70!black] at (-0.2, 0.3) {Objective: $z = 22 - 2x_3 - \tfrac12 x_4 \implies$ all coeff $\le 0$ (Stop)};

    % Connecting arrows between ledger cards
    \draw[->, thick, trapred] (4.15, 3.35) -- (4.15, 3.15);
    \draw[->, thick, trapred] (4.15, 1.65) -- (4.15, 1.45);
  \end{scope}
\end{tikzpicture}
\caption{The simplex walk in geometry and algebra. (a) Starting at the origin
$O(0,0)$ where both decision variables are nonbasic, Pivot~1 enters $x_2$ and
slides along edge $OC$ until constraint $x_3$ binds at $C(0,4)$, lifting profit
from $0$ to $20$. Pivot~2 enters $x_1$ and traverses edge $CB$ until $x_4$ binds
at $B(4,2)$, reaching the maximum profit $22$. Vertex $A$ is bypassed entirely.
(b) The corresponding dictionary ledger tracking entering/leaving variables,
ratio tests, and the termination condition.}
\label{ch28:fig-walk}
\end{figure}

\begin{labbox}
Trace the path on Figure~\ref{ch28:fig-walk}: simplex walked $O \to C \to B$, climbing
$0 \to 20 \to 22$. It never visited $A$. Now \emph{you} re-run it but break the
first tie differently: at iteration 0 let $x_1$ enter instead of $x_2$
(``largest coefficient'' is just one rule; ``first positive'' is another).
Do the ratio test, pivot, and keep going. You should travel $O \to A \to B$
and still land at $z=22$. The two paths reach the same optimum, as
convexity guarantees. Write out all three dictionaries yourself.
\end{labbox}

\subsection{Simplex as pseudocode}

The example contains the whole algorithm. Here it is in skeleton form,
operating on a slack-form dictionary with basic set $B$ and nonbasic set $N$.

\begin{algorithm}[H]
\caption{\textsc{Simplex} (one phase, on a feasible dictionary)}
\KwIn{slack-form dictionary: objective row $z = v + \sum_{j\in N} \bar c_j x_j$
  and rows $x_i = \bar b_i - \sum_{j\in N}\bar a_{ij}x_j$ for $i \in B$}
\KwOut{an optimal basic feasible solution, or report ``unbounded''}
\While{some $\bar c_e > 0$ for $e \in N$}{
  pick entering index $e \in N$ with $\bar c_e > 0$\;
  \tcp{ratio test: how far can $x_e$ rise?}
  \If{$\bar a_{ie} \le 0$ for all $i \in B$}{
    \Return ``unbounded'' \tcp*{$x_e$ rises forever, $z\to\infty$}
  }
  pick leaving index $\ell \in B$ minimizing $\bar b_i / \bar a_{ie}$ over
    $i$ with $\bar a_{ie} > 0$\;
  \tcp{pivot: swap $x_e$ into the basis, $x_\ell$ out}
  rewrite row $\ell$ to express $x_e$, then substitute into all other rows and
    the objective\;
  update $B \gets (B \setminus \{\ell\}) \cup \{e\}$ and
    $N \gets (N \setminus \{e\}) \cup \{\ell\}$\;
}
\Return the basic solution: $x_j = 0$ for $j\in N$, $x_i = \bar b_i$ for $i\in B$\;
\end{algorithm}

\begin{focusbox}
The loop makes three decisions:
\begin{enumerate}
\item \textbf{Optimality test.} If no nonbasic coefficient $\bar c_e$ is
  positive, stop -- you are optimal.
\item \textbf{Entering rule.} Among positive $\bar c_e$, pick one (largest
  coefficient, or Bland's rule, etc.).
\item \textbf{Ratio test / leaving rule.} The entering variable rises until the
  first basic variable hits 0; that one leaves. If \emph{nothing} stops it, the
  LP is unbounded.
\end{enumerate}
\end{focusbox}

\subsection{Two loose ends: starting and cycling}

We glossed over two things that deserve to be named.

\textbf{Getting an initial feasible vertex.} We were lucky: the origin was
feasible (all $b_i \ge 0$), so the slacks gave a ready basic feasible solution.
When some $b_i < 0$, the origin is infeasible and cannot serve as the starting basis.
The fix is a preliminary ``Phase~I'' LP: introduce an auxiliary variable that
measures total infeasibility and minimize it with simplex itself. If the
minimum is 0, you have found a genuine feasible vertex to start the real
(``Phase~II'') optimization from; if it is positive, the original LP is
infeasible. So simplex bootstraps itself.

\textbf{Cycling.} The ratio test can tie, producing a \emph{degenerate} pivot
that changes the basis but not the vertex or the objective (you move 0
distance). In rare crafted cases a sequence of degenerate pivots can loop
forever, returning to a basis it already saw. The cure is a tie-breaking
discipline called \emph{Bland's rule}: among allowed entering and leaving
candidates, always choose the one with the smallest index. It can be proved
that Bland's rule never cycles, so simplex always terminates.

\begin{historybox}
Leonid Kantorovich formulated linear programming in 1939 to optimize production, earning the 1975 Nobel Memorial Prize in Economic Sciences. George Dantzig independently developed the general linear programming model and the simplex method in 1947 while working on Air Force planning problems; the word ``programming'' was military planning lingo from that era. There is a famous tale that as a graduate student Dantzig arrived late to class, copied two ``homework'' problems off the board, and solved them -- only to learn they were open problems in statistics. The simplex method is in \textbf{CLRS Chapter 29}, specifically Section~29.3 (the algorithm), Section~29.4 (duality), and Section~29.5 (the initial basic feasible solution / Phase~I). CLRS develops the same pivot-on-a-slack-form machinery we just traced; read it right after this chapter and it will feel familiar. A notable fact, covered in CLRS's notes, is that simplex can take exponentially many steps on adversarial inputs (the Klee--Minty cube) yet is fast in practice (explained decades later by Spielman and Teng's smoothed analysis); the first \emph{provably} polynomial methods were Khachiyan's ellipsoid algorithm (1979) and Karmarkar's interior-point method (1984).
\end{historybox}

\begin{complexitybox}
\textbf{Simplex cost.} Each pivot is a Gaussian-elimination-style rewrite of an
$m$-row dictionary in $n$ variables, costing $O(mn)$ arithmetic per iteration
(or $O(m^2)$ in the revised, matrix form). The number of pivots is the wild
card: typically $O(m)$ to $O(m+n)$ in practice, giving excellent real-world
speed, but worst-case \emph{exponential} for every known deterministic pivot
rule. By contrast, interior-point methods solve an LP to accuracy $\eps$ in a
provably polynomial number of iterations -- roughly $O(\sqrt{n}\,\log(1/\eps))$
Newton steps, each a linear solve. In short: simplex is fast in practice and
ugly in the worst case; interior-point is polynomial always and competitive in
practice. Solvers ship both.
\end{complexitybox}

\section{Duality: Every LP Has a Shadow}
\label{ch28:duality}

Every linear program (the \emph{primal}) has a companion linear program
called its \emph{dual}, formed from the exact same coefficients rearranged.
The two are locked together so tightly that solving one solves the other, and
the dual reveals something the primal never explicitly asked for: the
\emph{value of each scarce resource}.

\subsection{Where does the dual come from?}

Go back to the furniture LP and ask a direct question: \emph{without} solving
it, can we prove an \emph{upper bound} on the best possible profit?
Take the labor constraint $2x_1 + 2x_2 \le 12$ and multiply it by $\tfrac52$:
\[
5x_1 + 5x_2 \le 30 .
\]
Since $x_1 \ge 0$, certainly $3x_1 \le 5x_1$, so the profit obeys
\[
z = 3x_1 + 5x_2 \le 5x_1 + 5x_2 \le 30 .
\]
We just proved, with one line of algebra, that no plan earns more than \$30 --
\emph{without ever solving the LP}. Can we do better? Try combining the two
constraints. Pick nonnegative multipliers $y_1$ (for wood) and $y_2$ (for
labor):
\[
y_1(x_1 + 2x_2) + y_2(2x_1 + 2x_2) \le 8y_1 + 12y_2 .
\]
The left side is $(y_1 + 2y_2)x_1 + (2y_1 + 2y_2)x_2$. If we choose the $y$'s so
that this dominates the objective coefficient-by-coefficient -- that is,
\[
y_1 + 2y_2 \ge 3 \quad(\text{coeff. of }x_1),\qquad
2y_1 + 2y_2 \ge 5 \quad(\text{coeff. of }x_2),
\]
then, because all $x_j \ge 0$,
\[
z = 3x_1 + 5x_2 \le (y_1+2y_2)x_1 + (2y_1+2y_2)x_2 \le 8y_1 + 12y_2 .
\]
So \emph{any} nonnegative $y_1, y_2$ satisfying those two inequalities gives an
upper bound $8y_1 + 12y_2$ on the profit. Naturally, we want the
\emph{tightest} such bound -- the smallest $8y_1 + 12y_2$. And that is itself a
linear program!

\begin{focusbox}
The \textbf{dual} of our furniture LP is:
\[
\begin{aligned}
\text{minimize}\quad   & 8y_1 + 12y_2 \\
\text{subject to}\quad & y_1 + 2y_2 \ge 3 \\
                       & 2y_1 + 2y_2 \ge 5 \\
                       & y_1, y_2 \ge 0.
\end{aligned}
\]
One dual variable per primal \emph{constraint}; one dual constraint per primal
\emph{variable}; the objective and right-hand side swap roles; max becomes min;
$\le$ becomes $\ge$. The dual is the search for the best provable upper bound.
\end{focusbox}

The mechanical recipe, for an LP in standard form (max $c^\top x$ s.t.
$Ax \le b$, $x\ge 0$):

\begin{center}
\begin{tabular}{ll}
\toprule
\textbf{Primal (max)} & \textbf{Dual (min)} \\
\midrule
maximize $c^\top x$              & minimize $b^\top y$ \\
$m$ constraints $Ax \le b$       & $m$ variables $y \ge 0$ \\
$n$ variables $x \ge 0$          & $n$ constraints $A^\top y \ge c$ \\
right-hand side $b$              & objective $b$ \\
objective $c$                    & right-hand side $c$ \\
\bottomrule
\end{tabular}
\end{center}

In matrix form: the dual of ``$\max\, c^\top x$ s.t. $Ax\le b,\ x\ge 0$'' is
``$\min\, b^\top y$ s.t. $A^\top y \ge c,\ y \ge 0$.'' Transpose the matrix,
swap $b$ and $c$, flip the inequalities, flip max/min. The construction is symmetric:
the dual of the dual is the primal again.

\subsection{Weak and strong duality}

\begin{theorem}[Weak duality]
If $x$ is any feasible solution of the primal (a max LP) and $y$ is any
feasible solution of the dual (a min LP), then $c^\top x \le b^\top y$. Every
primal value is below every dual value.
\end{theorem}

We essentially proved this above: the chain $c^\top x \le (A^\top y)^\top x =
y^\top (Ax) \le y^\top b$ uses only $x \ge 0$, $A^\top y \ge c$, $y \ge 0$, and
$Ax \le b$. It is just bookkeeping with inequalities, but it
sandwiches the unknown optimum from both sides.

\begin{intuitionbox}
Picture the primal values (feasible profits) as numbers piling up from below
and the dual values (provable upper bounds) coming down from above. Weak
duality says the two piles \emph{never cross}: every bound from above is at
least every achievable profit from below. The question that remains is whether
they actually \emph{meet}. If a profit and a bound are ever equal, that profit
must be optimal (nothing can beat its own upper bound) and that bound must be
tight. Strong duality states that whenever the primal has a finite optimum, the
two bounds meet with zero gap.
\end{intuitionbox}

Figure~\ref{ch28:fig-sandwich} pictures the squeeze. The achievable primal
profits cannot rise above the line of provable dual bounds, and at the shared
optimum the two sides touch at the single value $22$.

\begin{figure}[H]
\centering
\begin{tikzpicture}[scale=1.0]
  % a vertical value axis
  \draw[->,thick] (0,-0.3) -- (0,4.2) node[above] {value};
  \foreach \y/\lab in {0/0,1/,2/,3/,4/} { \draw (-0.07,\y) -- (0.07,\y); }
  % primal feasible values rise from below (blue)
  \draw[sourceblue,very thick] (1.3,0) -- (1.3,2.2);
  \node[sourceblue,below,align=center,font=\footnotesize] at (1.3,-0.1) {primal\\[-1pt](max, from below)};
  \draw[->,sourceblue,thick] (1.3,1.6) -- (1.3,2.1);
  % dual feasible values come down from above (orange)
  \draw[quotientorange,very thick] (4.0,4.0) -- (4.0,2.2);
  \node[quotientorange,below,align=center,font=\footnotesize] at (4.0,-0.1) {dual\\[-1pt](min, from above)};
  \draw[->,quotientorange,thick] (4.0,2.8) -- (4.0,2.3);
  % the meeting value
  \draw[dashed,trapred,thick] (0,2.2) -- (4.6,2.2);
  \node[trapred,right] at (4.6,2.2) {\small $z^\star = 22$};
  \node[left] at (0,2.2) {\small $22$};
  \node[left] at (0,0) {\small $0$};
\end{tikzpicture}
\caption{Weak and strong duality as a squeeze. Every feasible primal value
(blue, pushing up) stays at or below every feasible dual value (orange, pressing
down); weak duality forbids them from crossing. Strong duality says the gap
closes exactly at the common optimum $z^\star = 22$.}
\label{ch28:fig-sandwich}
\end{figure}

\begin{theorem}[Strong duality]
If the primal LP has a finite optimum, then so does the dual, and their optimal
values are \emph{equal}: $\max c^\top x = \min b^\top y$.
\end{theorem}

This is not obvious and does not hold for general (nonlinear) optimization. For
LPs it is a theorem -- in fact simplex proves it constructively: when simplex
halts, the negated objective coefficients on the slack variables turn out to be
exactly an optimal dual solution. Let us verify on the furniture problem. The
primal optimum was $22$. Solve the dual: minimize $8y_1 + 12y_2$ s.t.
$y_1 + 2y_2 \ge 3$, $2y_1+2y_2 \ge 5$, $y \ge 0$. The optimum sits where both
constraints are tight: $y_1 + 2y_2 = 3$ and $2y_1 + 2y_2 = 5$. Subtracting,
$y_1 = 2$; then $2 + 2y_2 = 3$ gives $y_2 = \tfrac12$. Dual value:
$8(2) + 12(\tfrac12) = 16 + 6 = 22$. The two optima are equal, as strong duality says.

\begin{trapbox}
A frequent confusion: students expect the dual optimum to be \emph{bigger} than
the primal because ``it's an upper bound.'' At the optimum they are
\emph{equal}; the dual is bigger only at \emph{suboptimal} feasible points. Weak
duality gives the inequality everywhere; strong duality collapses it to
equality precisely at the joint optimum. Also: strong duality can fail if the
primal is infeasible or unbounded -- then the meeting-in-the-middle story needs
the careful case analysis CLRS Section~29.4 provides.
\end{trapbox}

\subsection{What the dual variables mean: shadow prices}

The dual values $y_1 = 2$ and $y_2 = \tfrac12$ have a concrete economic
meaning: they are \emph{shadow prices}.
Specifically, $y_i$ represents the rate at which the optimal profit would
increase if you were granted one additional unit of resource $i$.

Check it. We had 8 board-feet of wood; $y_1 = 2$ predicts that a 9th board-foot
is worth \$2 of extra profit. Indeed, re-solving with the wood limit raised to
9 moves the optimum to $(3,3)$ with $z = 24$ -- up by \$2, exactly $y_1$. The
labor shadow price $y_2 = \tfrac12$ says an extra hour of labor is worth fifty
cents at the margin.

\begin{aestheticbox}
Shadow prices make LP especially useful to managers. Solve the plan once and the dual
\emph{automatically} tells you what each bottleneck resource is worth -- where
to invest, what overtime is worth paying for, which contract to renegotiate.
A resource with a \emph{zero} shadow price has slack to spare; buying more of it
is wasted money. Along with a plan, the LP gives you the
\emph{prices of your constraints}. This is the same idea that underlies market
prices in competitive-equilibrium economics, and it is why Dantzig's method and
Koopmans's economics are closely related (Kantorovich and Koopmans won the 1975
Nobel in economics for related work).
\end{aestheticbox}

\begin{formalbox}
\textbf{Complementary slackness.}
At a primal-dual optimum, for each $i$ and $j$:
\[
y_i > 0 \;\Rightarrow\; \text{(primal constraint $i$ is tight)},\qquad
x_j > 0 \;\Rightarrow\; \text{(dual constraint $j$ is tight)}.
\]
Equivalently: a resource with leftover slack ($Ax < b$ in row $i$) has shadow
price $y_i = 0$; and you produce a positive amount of product $j$ ($x_j > 0$)
only if its dual constraint binds (it ``pays for itself''). In the furniture
optimum both resources were fully used, so both $y_i > 0$; and we built both
products, so both dual constraints were tight -- consistent with the solution we
found.
\end{formalbox}

\section{LP Relaxation as a Design Tool}
\label{ch28:relaxation}

We close with the reason an algorithms book cares about LP beyond ``it solves
LPs.'' Linear programming is a \emph{design tool} for attacking problems that
are not themselves LPs -- especially hard combinatorial ones.

The trouble usually is integrality. Real decisions are often whole numbers or
yes/no: you ship a truck or you don't; you assign a worker to a shift or you
don't; you pick a vertex into a cover or you leave it out. Add the constraint
``$x_j \in \{0,1\}$'' (or ``$x_j \in \Z$'') to an LP and you get an
\emph{integer linear program} (ILP). This one extra requirement makes the problem much
harder: ILP is NP-hard in general (Chapter~33), and the corner-walking
of simplex no longer finds integer corners.

The standard workaround is \emph{LP relaxation}: drop the integrality requirement.
Replace ``$x_j \in \{0,1\}$'' with the interval ``$0 \le x_j \le 1$,'' solve
the resulting ordinary LP (fast!), and use its answer as a guide.

\begin{intuitionbox}
The relaxed LP is more permissive -- it allows fractional, ``half a truck''
solutions that the integer problem forbids. So the relaxed feasible region
\emph{contains} all the integer points and then some. For a minimization
problem, allowing more options can only lower (or equal) the best value, so
\[
\text{LP optimum} \;\le\; \text{ILP optimum}.
\]
The relaxation gives you a \emph{bound} for free, and a fractional solution to
round. That bound drives two important algorithmic
techniques in the next chapters.
\end{intuitionbox}

Figure~\ref{ch28:fig-relax} makes the containment concrete. The shaded LP
polytope holds every legal integer point (the dots) but also a continuum of
fractional points between them; the LP optimum can sit at a fractional corner
the ILP could never reach.

\begin{figure}[H]
\centering
\begin{tikzpicture}[scale=0.85]
  \draw[->,thick] (-0.3,0) -- (5,0) node[right] {$x_1$};
  \draw[->,thick] (0,-0.3) -- (0,4.2) node[above] {$x_2$};
  % a polytope containing some lattice points
  \fill[gray!18,opacity=0.8] (0,0) -- (4,0) -- (3.5,3) -- (0,2.5) -- cycle;
  \draw[structurepurple,thick] (0,0) -- (4,0) -- (3.5,3) -- (0,2.5) -- cycle;
  % integer lattice points inside (filled) and the fractional LP corner (open star)
  \foreach \x in {0,1,2,3,4} \foreach \y in {0,1,2,3} {
    \fill[black!55] (\x,\y) circle (0.8pt);
  }
  \foreach \p in {(0,0),(1,0),(2,0),(3,0),(4,0),(0,1),(1,1),(2,1),(3,1),
                  (0,2),(1,2),(2,2),(3,2)} {
    \fill[targetgreen] \p circle (2.0pt);
  }
  \fill[trapred] (3.5,3) circle (2.4pt);
  \node[trapred,right] at (3.5,3) {\small LP corner (fractional)};
  \node[targetgreen,below right] at (4,0) {};
  \node[targetgreen] at (2.3,-0.9) {\small green dots: feasible integer points};
\end{tikzpicture}
\caption{LP relaxation. The relaxed polytope (shaded, purple boundary) contains
all the integer feasible points (green dots) and the fractional points between
them. The relaxed optimum may land on a fractional corner (red), which is why
its value bounds the integer optimum but its coordinates may need rounding.}
\label{ch28:fig-relax}
\end{figure}

\textbf{Use 1: bounding for exact search.} Branch-and-bound solves ILPs exactly
by exploring a tree of subproblems. At each node it solves the LP relaxation; if
the relaxed optimum is already worse than the best integer solution found so
far, the entire subtree can be pruned -- no integer point down there can beat
what we have. The relaxation's bound does the pruning. This is how industrial
ILP solvers (CPLEX, Gurobi) solve problems with millions of variables.

\textbf{Use 2: rounding for approximation.} When exact is too slow, solve the
relaxation, then \emph{round} the fractional solution to a nearby integer one,
and prove the rounding doesn't cost too much. Consider \textsc{Vertex Cover}:
choose a smallest set of vertices touching every edge. The ILP has $x_v \in
\{0,1\}$ and a constraint $x_u + x_v \ge 1$ for each edge $(u,v)$. Relax to
$0 \le x_v \le 1$, solve the LP, and round: put $v$ in the cover iff
$x_v \ge \tfrac12$.

\begin{focusbox}
Why the rounding is safe and good: every edge constraint $x_u + x_v \ge 1$
forces at least one of $x_u, x_v$ to be $\ge \tfrac12$, so the rounded set still
covers every edge (still feasible). And we at most \emph{double} each variable
(from $\ge\tfrac12$ up to $1$), so the rounded cost is at most
$2 \times (\text{LP optimum}) \le 2 \times (\text{ILP optimum})$. That is a
\textbf{2-approximation}: a polynomial-time algorithm guaranteed within a
factor of 2 of optimal, with the LP relaxation supplying both the rounding
recipe and the proof. We meet this style again in Chapter~34.
\end{focusbox}

\begin{walkbox}
Take this on a walk. The shortest-path-as-an-LP from
Section~\ref{ch28:modeling} and the maximum-flow problem (which can also be
written as an LP) both have a special structural property: when capacities or
weights are integers, their LP relaxations \emph{automatically} have integer
optimal vertices -- you never
need to round. Why should some combinatorial LPs hand you whole numbers for
free while \textsc{Vertex Cover} does not? Think about what is special about the
constraint matrices of flow and shortest-path problems. (The keyword, for
after your walk, is \emph{total unimodularity}: matrices whose square submatrices
all have determinant $0$ or $\pm1$ guarantee integer vertices. Network-flow
matrices are totally unimodular; the vertex-cover matrix is not.)
\end{walkbox}

\begin{focusbox}
LP relaxation links continuous optimization to combinatorial problems.
It connects directly to integer programming, NP-completeness (Chapter~33, CLRS
Chapter~34), and approximation algorithms (Chapter~34, CLRS Chapter~35), where
LP-rounding and primal-dual methods are central design strategies. CLRS
Chapter~29 covers pure LP: standard and slack forms (29.1), formulating problems
(29.2), simplex (29.3), duality (29.4), and the initial solution (29.5).
\end{focusbox}

\section{Summary}
\label{ch28:summary}

Starting from a two-variable workshop problem, we developed the core machinery
of linear optimization:

\begin{itemize}
\item An \textbf{LP} maximizes a linear objective subject to linear
  constraints. Modeling means naming decision variables and writing the
  objective and constraints in their terms -- staying strictly linear.
\item \textbf{Standard form} (max, $\le$, $x \ge 0$) is the canonical shape;
  \textbf{slack form} turns every inequality into an equality with a slack
  variable, splitting variables into basic and nonbasic.
\item \textbf{Geometry}: the feasible region is a convex polytope, the optimum
  lives at a vertex, and you only ever need to check corners.
\item The \textbf{simplex method} walks from vertex to adjacent vertex, always
  uphill, via entering/leaving pivots, stopping when no objective coefficient is
  positive. Fast in practice, exponential in the worst case.
\item \textbf{Duality} pairs every LP with a shadow LP. Weak duality sandwiches
  the optimum; strong duality says primal and dual optima are equal; the dual
  variables are shadow prices that value each scarce resource.
\item \textbf{LP relaxation} drops integrality to turn hard integer problems
  into solvable LPs, giving bounds for branch-and-bound and rounding recipes for
  approximation, which leads into the book's final chapters.
\end{itemize}

Linear programming optimizes linear forms over flat geometric polytopes. In the next chapter, we turn from straight boundaries to curves: polynomials, their dual coefficient and point-value representations, and the fast Fourier transform that multiplies degree-$n$ polynomials in near-linear time.

\section*{Exercises}
\addcontentsline{toc}{section}{Exercises}

\begin{exercise}[Model it]
A bakery makes muffins and cookies. A muffin earns \$2 and uses 3 oz flour and
1 egg; a cookie earns \$1 and uses 1 oz flour and 1 egg. There are 30 oz of
flour and 12 eggs. Write the LP that maximizes revenue. Identify the decision
variables, objective, and constraints.
\end{exercise}

\begin{exercise}[Corners by hand]
For the bakery LP above, draw the feasible region in the muffin--cookie plane,
list all its vertices, evaluate the revenue at each, and read off the optimal
production plan. Confirm the optimum is at a corner.
\end{exercise}

\begin{exercise}[Into standard form]
Convert to standard form: ``minimize $4x_1 - x_2 + 2x_3$ subject to
$x_1 - x_2 \ge 3$, $x_2 + x_3 = 5$, $x_1 \ge 0$, $x_2 \ge 0$, $x_3$ free.''
Show every transformation step.
\end{exercise}

\begin{exercise}[Slack form and the origin]
Put the furniture LP's constraints into slack form, then write the basic
feasible solution you get by setting both original variables to 0. Which
constraints are tight there? Now do the same after raising the wood limit from
8 to 10 -- does the origin remain feasible?
\end{exercise}

\begin{exercise}[Run simplex]
Solve by simplex, showing every dictionary:
``maximize $x_1 + 2x_2$ subject to $x_1 + x_2 \le 4$, $x_1 + 3x_2 \le 6$,
$x_1, x_2 \ge 0$.'' Report the optimal vertex, the optimal value, and the path
of vertices simplex visited.
\end{exercise}

\begin{exercise}[Spot the unbounded LP]
Try to run simplex on ``maximize $2x_1 + x_2$ subject to $x_1 - x_2 \le 1$,
$x_1, x_2 \ge 0$.'' At some iteration the ratio test will find no positive
coefficient in the entering column. Explain why, and explain geometrically why
this LP has no finite maximum.
\end{exercise}

\begin{exercise}[Write the dual]
Write the dual of the bakery LP from Exercise~28.1. Identify what each dual
variable would mean economically (the shadow price of which resource?).
\end{exercise}

\begin{exercise}[Verify strong duality]
Solve the dual of the furniture LP (minimize $8y_1 + 12y_2$ subject to
$y_1 + 2y_2 \ge 3$, $2y_1 + 2y_2 \ge 5$, $y \ge 0$) by checking its vertices.
Confirm the optimal dual value equals the primal optimum of 22, and state which
primal constraints the nonzero $y_i$ correspond to.
\end{exercise}

\begin{exercise}[Shadow price prediction]
Using the furniture shadow prices $y_1 = 2$ (wood) and $y_2 = \tfrac12$
(labor), predict the new optimal profit if labor rises from 12 to 14 hours.
Then re-solve the LP from scratch and check your prediction. For roughly how
large an increase do you expect the shadow price to stay valid? (Hint: think
about when the optimal corner changes.)
\end{exercise}

\begin{exercise}[Relax and round]
Take the triangle graph with vertices $\{a,b,c\}$ and edges $ab, bc, ca$. Write
the \textsc{Vertex Cover} ILP, then its LP relaxation. Show that
$x_a = x_b = x_c = \tfrac12$ is an optimal \emph{fractional} solution with value
$\tfrac32$. Round by the ``$x_v \ge \tfrac12$'' rule and report the resulting
integer cover. Compare its size to the true optimum (which is 2) and to the
$2 \times \tfrac32 = 3$ bound the analysis promises.
\end{exercise}

\begin{exercise}[Challenge: shortest path is an LP]
Take a 4-vertex graph $s \to a \to t$ and $s \to t$ with weights
$w(s,a)=2$, $w(a,t)=3$, $w(s,t)=6$. Write the shortest-path LP from
Section~\ref{ch28:modeling} (maximize $\sum_v d_v$ subject to
$d_v - d_u \le w(u,v)$ and $d_s = 0$). Solve it and check that the $d_v$ you
get are the true shortest-path distances. Why does \emph{maximizing} produce the
\emph{shortest} distances?
\end{exercise}

\begin{exercise}[Challenge: dual of the dual]
Starting from the standard-form dual ``minimize $b^\top y$ subject to
$A^\top y \ge c$, $y \ge 0$,'' first convert it into a maximization in standard
form, then take \emph{its} dual using the recipe in Section~\ref{ch28:duality}.
Show you recover the original primal ``maximize $c^\top x$ subject to
$Ax \le b$, $x \ge 0$.'' Conclude that dualizing twice is the identity.
\end{exercise}
